% OC Core 1.3 — Cross-level structures master module
% Diese Datei bündelt alle Cross-K-Verknüpfungsabschnitte.

\section{Cross-Level-Strukturen}
\label{sec:crossk_master}

\subsection{Zweck der Cross-K-Ebene}
Die Cross-K-Ebene formalisiert alle Strukturen, Übergänge und Beschränkungen, die
\emph{verschiedene Kontinuumsstufen $K_0 \rightarrow K_{12}$ verknüpfen}.
Während jede $K$-Stufe eine eigene Ontologie besitzt (Achsen, Potentiale,
Schwellen, Flüsse, Zyklen, Domänen $\Omega(K)$, Grenzen $\partial\Omega(K)$ und
das Kontinuümsmaß $k(K)$), erfordern viele Phänomene im Core
\textbf{Relationen zwischen Stufen}. Diese Relationen sind weder optional
noch extern: Sie gehören zur intrinsischen Architektur eines dynamischen
Kontinuums.

Die Cross-K-Ebene stellt damit:
\begin{itemize}
  \item \textbf{Abbildungsstrukturen} zwischen Stufen bereit,
  \item \textbf{Übergangsoperatoren} und ihre Nebenbedingungen spezifiziert,
  \item \textbf{Konsistenzbedingungen} einführt, die die Entwicklung
        über benachbarte Stufen hinweg mathematisch und konzeptuell korrekt
        machen,
  \item \textbf{globale Landschaftsregeln} vorgibt, die beschreiben, wie der
        gesamte Stapel der $K$-Stufen ein kohärentes Meta-Kontinuum bildet.
\end{itemize}

\subsection{Definition von Cross-K-Strukturen}
Eine \emph{Cross-K-Struktur} ist jedes formale Objekt, das:
\begin{enumerate}
  \item mindestens zwei Stufen $K_i, K_j$ mit $i \neq j$ umfasst;
  \item eine Relation $R_{ij}$ zwischen ihren Komponenten
        $(A,P,\Theta,J,C,\Omega,\partial\Omega,k)$ festlegt;
  \item die allgemeine Kompatibilitätsregel erfüllt:
        \[
        R_{ij} : K_i \leftrightarrow K_j
        \]
        unter Erhaltung der Kontinuität und erlaubten Übergängen im Rahmen des
        globalen Operators $E$.
\end{enumerate}

Damit umfassen Cross-K-Strukturen:
\begin{itemize}
  \item \textbf{Übergangsabbildungen} (dimensional, strukturell, energetisch),
  \item \textbf{Erbbeziehungen} (Achsen, Potentiale, Schwellen),
  \item \textbf{Randtransformationen} zwischen $\partial\Omega(K_i)$ und
        $\partial\Omega(K_j)$,
  \item \textbf{Zyklische Projektionen und Hebungen} (Abbildung von Zyklen
        zwischen Stufen),
  \item \textbf{globale Nebenbedingungen}, die sicherstellen, dass keine Stufe
        die Regeln ihrer Nachbarn verletzt.
\end{itemize}

\subsection{Rolle dieser Master-Datei}
Dieses Master-Modul ist bewusst eine \textbf{dünne Orchestrierungsebene}.
Es enthält keine umfangreichen theoretischen Inhalte. Stattdessen:
\begin{itemize}
  \item führt das Konzept der Cross-K-Strukturen ein,
  \item stellt die globale Definition und den Zweck dar,
  \item und führt alle detaillierten Stufen-analysen über \verb|\input| zusammen.
\end{itemize}

Die einzelnen Übergangskapitel sind folgendermaßen organisiert:
\begin{itemize}
  \item globale Landschaft (Alle-Stufen-Perspektive),
  \item lokale Übergänge ($K_i \rightarrow K_{i+1}$),
  \item besondere bidirektionale oder nicht benachbarte Abbildungen, wenn nötig.
\end{itemize}

\subsection{Modulstruktur}
\noindent Folgende Dateien werden von diesem Modul orchestriert:
\begin{itemize}
  \item \verb|crossk_global_landscape.tex| — globale Mehrstufen-Geometrie;
  \item \verb|crossk_k0_k1.tex|, \verb|crossk_k1_k2.tex|, \dots — Übergänge
        benachbarter Stufen;
  \item \verb|crossk_k10_k11.tex|, \verb|crossk_k11_k12.tex| — Übergänge der
        oberen Stufen und Meta-Strukturen.
\end{itemize}

\subsection{Einbindung der Detailkapitel}
\noindent Die detaillierten Cross-K-Analysen können wie folgt eingebunden werden:
\begin{verbatim}
\input{content/crossk/crossk_global_landscape.tex}
\input{content/crossk/crossk_k0_k1.tex}
\input{content/crossk/crossk_k1_k2.tex}
...
\input{content/crossk/crossk_k11_k12.tex}
\end{verbatim}

\noindent Diese Datei definiert somit den \textbf{konzeptuellen Container} für den
Cross-K-Bereich von OC Core 1.3.
